% a04_Bh_bulk2  Proof of app:a04:lem:B2, part 2: the line lemmas for row (III) (coarse lines, free lines with differences, outer lines, power counting)

\subsubsection*{B2.5 The lines of row $(\mathrm{III})$}

\begin{lemma}[coarse lines at spacing $\tau$]\label{app:a04:lem:Gam}
Let $\theta\in(0,1)$. There is $C=C(\theta,\mathfrak a,m_0)$ such that for all $\tau\in(0,1]$ ($\tau\le\frac15$ for $f_D$), $m'\in(0,m_0]$, $b\in[\mathfrak a(1-5^{-2}),\mathfrak a]$, $x,x',h,h'\in\Lambda_\tau$ and every forward difference $\nabla=\nabla^{(\tau)}_\mu$ acting on $x$ or on $x'$:
\begin{enumerate}[label=\textup{(\ensuremath{\Gamma}\arabic*)},leftmargin=*]
\item $0\le\Gamma_\tau(x,x)\le C$, $|\Gamma_\tau(x,x')|\le C$;
\item $|\nabla\Gamma_\tau(x,x')|\le C$;
\item $|\nabla\Gamma_\tau(x+h,x'+h')-\nabla\Gamma_\tau(x,x')|\le C(|h|+|h'|)^\theta$;
\item for $x,x'\in\Lambda_t$ ($t=5^{-j}$, both members of the diagonal pair, each with its own weight): $|\Gamma_{t/5}(x,x')-\Gamma_t(x,x')|\le Ct$ and $|\nabla^{(t/5)}\Gamma_{t/5}(x,x')-\nabla^{(t)}\Gamma_t(x,x')|\le Ct(1+\ln t^{-1})^3$;
\item $|\partial_b\Gamma_\tau(x,x')|\le C$, $|\partial_b\nabla\Gamma_\tau(x,x')|\le C$.
\end{enumerate}
The same holds for the photon coarse line $\Gamma^\ph_\tau$ with $m'\to\mu'$ (uniform weight).
\end{lemma}
\begin{proof}
In the Fourier representation a forward difference in $x$ multiplies the $n$-th term of $\hat v^\tau_x$ by $d_\tau(p_\mu)$, $|d_\tau(\vartheta)|\le|\vartheta|$, and a displacement $x\to x+h$ by $e^{-ip\cdot h}$, $|e^{-ip\cdot h}-1|\le\min(2,|p||h|)\le2^{1-\theta}|p|^\theta|h|^\theta$. Write $(*)_\kappa$ for \eqref{app:a04:eq:star}, resp.\ (A2) for $f_D$, and use $\Mh\le C(P^2+m'^2)$ ((A3) for $f_D$).
($\Gamma$1) is \cref{app:a04:lem:B0}(i), including its Dirichlet-weight paragraph.
($\Gamma$2): $|\nabla\hat v_x|\le(*)_1$, $|\hat v_{x'}|\le(*)_0$: the integrand is $\le C(P^2+m'^2)^{-1/2}$, integrable on $\mathbb T$ uniformly in $m'$ (also at $m'=0$).
($\Gamma$3): displacing $x$: $|\sum_na_\tau d_\tau(p_\mu)(e^{-ip\cdot h}-1)e^{-ip\cdot x}|\le C|h|^\theta(*)_{1+\theta}\le C|h|^\theta(P^2+m'^2)^{(\theta-1)/2}$, times $|\hat v_{x'}|\Mh\le C$; displacing $x'$: $|\hat v_{x'+h'}-\hat v_{x'}|\le C|h'|^\theta(*)_\theta\le C|h'|^\theta(P^2+m'^2)^{(\theta-2)/2}$, times $|\nabla\hat v_x|\Mh\le C(P^2+m'^2)^{1/2}$; in both cases the integrand is $C|h|^\theta(P^2+m'^2)^{(\theta-1)/2}$, integrable since $(\theta-1)/2>-3/2$. For $f_D$ the D-alias part of $(*)_{1+\theta}$ is $C\tau^{1-\theta}(1+\ln\tau^{-1})^3\le C_\theta$ (A2), so the constant stays $\tau$-uniform. ($\theta=1$ fails only because $S_2=\infty$.)
($\Gamma$4), first clause: \cref{app:a04:lem:B0}(ii) at both spacings, each member compared with the same $\Gamma_0$ ($\Lambda_t\subset\Lambda_{t/5}$; for $f_D$ by (A4), each with its own weight): $\le Ct/5+Ct$. Second clause: compare both with $\partial\Gamma_0(x,x')\defeq\int\overline{\sum_na_0(n)(-ip_\mu)e^{-ip\cdot x}}\,\hat v^0_{x'}\Mh_0\dP$. Telescoping $\nabla\hat v^\tau_x\,\hat v^\tau_{x'}\Mh_\tau-\partial\hat v^0_x\,\hat v^0_{x'}\Mh_0$, the first term is $|\nabla\hat v^\tau_x-\partial\hat v^0_x|\,|\hat v^\tau_{x'}|\Mh_\tau\le C\tau(1+\ln\tau^{-1})^3$ (the derivative clause of (A4); for the uniform weight the same computation without D-aliases), and the other two, $|\partial\hat v^0_x||\hat v^\tau_{x'}-\hat v^0_{x'}|\Mh_\tau$ and $|\partial\hat v^0_x||\hat v^0_{x'}||\Mh_\tau-\Mh_0|$, are $\le C\tau(P^2+m'^2)^{-1/2}\cdot C$ by \S B2.3; integrate over $\mathbb T$ and apply at $\tau=t$ and $\tau=t/5$.
($\Gamma$5): $\partial_b\Mh=b^{-2}\Mh^2\le C(P^2+m'^2)^2$, so $|\partial_b\Gamma_\tau|\le C\int(P^2+m'^2)^{-2+2}\dP\le C$ and $|\partial_b\nabla\Gamma_\tau|\le C\int(P^2+m'^2)^{-1/2-1+2}\dP\le C$.
Photon: the same formulas with $m'\to\mu'$, the same $Q$ and $b$ \cite[(2.16)--(2.17)]{K1}.
\end{proof}

\begin{lemma}[the free lattice propagator with differences]\label{app:a04:lem:LPgrad}
For forward lattice differences $\nabla^\alpha_t$ of order $|\alpha|\le2$ there is $C$ such that for all $t>0$, $m'\ge0$ and $r\in t\Z^3\setminus0$,
\[
|\nabla^\alpha_tC_t(r;m')-\partial^\alpha C_0(r;m')|\le C\,t\,|r|^{-2-|\alpha|}.
\]
The same holds for the photon kernel $D_t(\cdot;\mu')$.
\end{lemma}
\begin{proof}
In the notation of the proof of \cref{app:a04:lem:LP}, for $r=tn$: $\nabla^\alpha_tC_t(tn)=t^{-1-|\alpha|}\int_0^\infty e^{-M^2\tau}(\nabla^\alpha_1P_\tau)(n)\,d\tau$ and $\partial^\alpha C_0(tn)=t^{-1-|\alpha|}\int_0^\infty e^{-M^2\tau}(\partial^\alpha G_\tau)(n)\,d\tau$, with $\nabla_1$ the unit-lattice difference, so it suffices to show $\int_0^\infty|\nabla^\alpha_1P_\tau-\partial^\alpha G_\tau|(n)\,d\tau\le C|n|^{-2-|\alpha|}$ ($t^{-1-|\alpha|}|n|^{-2-|\alpha|}=t|r|^{-2-|\alpha|}$). \emph{One dimension.} $\nabla^jq_\tau(\ell)=(2\pi)^{-1}\int_{-\pi}^\pi e^{i\ell\vartheta}(e^{i\vartheta}-1)^je^{-2\tau(1-\cos\vartheta)}d\vartheta$ and $\partial^jg_\tau(\ell)=(2\pi)^{-1}\int_\R e^{i\ell\vartheta}(i\vartheta)^je^{-\tau\vartheta^2}d\vartheta$. Repeat the one-dimensional argument of \cref{app:a04:lem:LP} with these insertions: on the shifted contour $|e^{iz}-1|\le C|z|$ and $|(e^{iz}-1)^j-(iz)^j|\le C|z|^{j+1}$. The term $(iz)^j(e^{-\tau\varphi}-e^{-\tau\psi})$ carries one more factor $|z|^j\le C(\vartheta^2+\xi^2)^{j/2}$, which turns $\tau^{-3/2}$ into $\tau^{-(3+j)/2}$; the term $[(e^{iz}-1)^j-(iz)^j]e^{-\tau\varphi}$ is $O(|z|^{j+1})$ and gives $\tau^{-(2+j)/2}$. Hence, for $\tau\ge1$, $|\nabla^jq_\tau(\ell)-\partial^jg_\tau(\ell)|\le C\tau^{-(2+j)/2}e^{-c\min(\ell^2/\tau,|\ell|)}$, and separately $|\nabla^jq_\tau(\ell)|,|\partial^jg_\tau(\ell)|\le C\tau^{-(1+j)/2}e^{-c\min(\ell^2/\tau,|\ell|)}$ for all $\tau>0$ (the same contour shift). \emph{Three dimensions.} $\nabla^\alpha_1P_\tau(n)=\prod_\mu\nabla^{\alpha_\mu}q_\tau(n_\mu)$; telescope as in \cref{app:a04:lem:LP} with $\varrho=|n|_\infty\ge1$. For $\tau\ge1$ the integrand is $\le C\tau^{-(4+|\alpha|)/2}e^{-c\min(\varrho^2/\tau,\varrho)}$ (one differenced factor at $\tau^{-(2+j)/2}$, the two others at their own decay); $\int_1^\varrho\le Ce^{-c\varrho}$ and $\int_\varrho^\infty\tau^{-(4+|\alpha|)/2}e^{-c\varrho^2/\tau}d\tau\le C\varrho^{-2-|\alpha|}$. For $\tau\le1$ write each forward difference as a sum of at most $2^{|\alpha|}$ shifted values: $|\nabla^\alpha_1P_\tau(n)|\le2^{|\alpha|}\max_{|n'-n|_\infty\le2}|P_\tau(n')|\le Ce^{-c(\varrho-2)}$, since the contour shift at $\xi=\operatorname{sgn}(\ell)$ gives $|q_\tau(\ell)|\le e^{-|\ell|+2\tau(\cosh1-1)}\le Ce^{-|\ell|}$ for $\tau\le1$ (the shifted $\vartheta$-integrand has modulus $\le1$); and $|\partial^\alpha G_\tau(n)|\le C\tau^{-(3+|\alpha|)/2}e^{-\varrho^2/4\tau}$, integrable with $\int_0^1\le Ce^{-\varrho^2/8}$. The finitely many $n$ with $\varrho\le3$ are covered by $\int_0^\infty(1+\tau)^{-(3+|\alpha|)/2}d\tau<\infty$ for the lattice term (bounded for $\tau\le1$ by the same shifted-value argument, $\le C\tau^{-(3+|\alpha|)/2}$ for $\tau\ge1$ by the separate bounds) and by the Gaussian factor for the continuum term. The mass enters only through $0<e^{-M^2\tau}\le1$.
\end{proof}

\begin{lemma}[free and outer lines]\label{app:a04:lem:TIOUT}
Let $w_s(r)\defeq(|r|+\tau)^{-s}$ on $\Lambda_\tau$. Uniformly in $\tau\in(0,1]$ and $m',\mu'\ge0$:
\begin{enumerate}[label=\textup{(T\arabic*)},leftmargin=*]
\item $|\nabla^\alpha C_\tau(r)|\le Cw_{1+|\alpha|}(r)$ for all $r\in\Lambda_\tau$ and $|\alpha|\le2$; the same for $D_\tau$.
\item For $r\in\Lambda_t\setminus0$: $|C_{t/5}(r)-C_t(r)|\le Ct^2|r|^{-3}\le Ctw_2(r)$ and $|\nabla^{(t/5)}C_{t/5}(r)-\nabla^{(t)}C_t(r)|\le Ct|r|^{-3}\le Ctw_3(r)$.
\item For $r,r'\in\Lambda_\tau$ with $|r|\ge8t$, $|r'-r|\le2\sqrt3t$, $\tau\le t$ and $|\alpha|\le1$: $|\nabla^\alpha C_\tau(r')-\nabla^\alpha C_\tau(r)|\le Ctw_{2+|\alpha|}(r)$.
\end{enumerate}
For the outer line $G_\tau=C_\tau-\Gamma_\tau$ (both weights; $b\in[\mathfrak a(1-5^{-2}),\mathfrak a]$) and $i\le1$:
\begin{enumerate}[label=\textup{(O\arabic*)},leftmargin=*]
\item $|\nabla^iG_\tau(x,y)|\le Cw_{1+i}(x-y)$ for $|x-y|\le1$ and $|\nabla^iG_\tau(x,y)|\le C$ for all $x,y$; hence on a ball $|x-y|\le R$, $|\nabla^iG_\tau(x,y)|\le C(R+1)^{1+i}w_{1+i}(x-y)$.
\item With $c\defeq\alpha_0/(2\sqrt3)$ and $R_\tau\defeq(16\sqrt3/\alpha_0)\ln\tau^{-1}+1$: $|\nabla^iG_\tau(x,y)|\le C\tau^4e^{-c|x-y|}$ for $|x-y|\ge R_\tau$, where $\alpha_0=\alpha_0(\mathfrak a,\norm{f}_\infty)>0$ is the constant of \cref{app:a04:lem:CT}, independent of $\tau$, $m'$ and $b$.
\end{enumerate}
The photon lines $D_\tau$, $G^\ph_\tau$ obey the same bounds with $\mu'$.
\end{lemma}
\begin{proof}
(T1): by \cref{app:a04:lem:LPgrad}, $|\nabla^\alpha C_\tau(r)-\partial^\alpha C_0(r)|\le C\tau|r|^{-2-|\alpha|}\le C|r|^{-1-|\alpha|}$ for $r\ne0$ ($\tau\le|r|$), and $|\partial^\alpha C_0(r)|\le C|r|^{-1-|\alpha|}$ for the Yukawa kernel $C_0=e^{-m'|r|}/(4\pi|r|)$, uniformly in $m'\ge0$; at $r=0$, $C_\tau(0)=\int_{\mathrm{BZ}_\tau}\hat C_\tau\dP\le C\tau^{-1}$ and $|\nabla^\alpha C_\tau|\le2^{|\alpha|}\tau^{-|\alpha|}C_\tau(0)$ ($|C_\tau(r)|\le C_\tau(0)$ by positive definiteness). Together $|\nabla^\alpha C_\tau(r)|\le C(|r|+\tau)^{-1-|\alpha|}$. (T2): \cref{app:a04:lem:LP} at both spacings, $|C_{t/5}-C_0|+|C_0-C_t|\le2C_*t^2|r|^{-3}$, and $t^2|r|^{-3}\le4tw_2(r)$ since $|r|\ge t$; \cref{app:a04:lem:LPgrad} likewise for the differences (the coarse forward difference of $C_{t/5}$ at spacing $t/5$ is taken at the points of $\Lambda_t$ with step $t/5$, as the statement says). (T3): a lattice path in $\Lambda_\tau$ from $r$ to $r'$ has $\le3\cdot2\sqrt3t/\tau$ steps; along it $|r''|\ge|r|-2\sqrt3t\ge|r|/2$; each step changes $\nabla^\alpha C_\tau$ by $\le\tau\sup|\nabla^{\alpha+e_\nu}C_\tau(r'')|\le C\tau|r|^{-2-|\alpha|}2^{2+|\alpha|}$ by (T1) at order $|\alpha|+1\le2$; the sum is $\le Ctw_{2+|\alpha|}(r)$. (O1): (T1) and ($\Gamma$1), ($\Gamma$2) of \cref{app:a04:lem:Gam}; on $|r|\le R$, $C\le C(R+1)^sw_s(r)$. (O2): \cref{app:a04:lem:CT} at unit block scale ($\eta=1$, $H=-\Delta_\tau+m'^2+bM_f\mathsf PM_f$, $b\in[\mathfrak a(1-5^{-2}),\mathfrak a]$, block averages of $f$ equal to $1$): $|G_\tau(x,y)|\le2c_0^{-1}\tau^{-3}e^{-\alpha_0|x-y|_\infty}$, so $|\nabla G_\tau|\le C\tau^{-4}e^{-\alpha_0(|x-y|_\infty-\tau)}$; with $|r|_\infty\ge|r|/\sqrt3$, $e^{-\alpha_0|r|/\sqrt3}=e^{-c|r|}e^{-c|r|}\le e^{-cR_\tau}e^{-c|r|}\le\tau^8e^{-c|r|}$ for $|r|\ge R_\tau$. Photon: (T1)--(T3) with $\mu'$ (\cref{app:a04:lem:LP,app:a04:lem:LPgrad} hold for every mass), (O1)--(O2) with the uniform weight.
\end{proof}

\begin{lemma}[power counting on confined triangles]\label{app:a04:lem:PC}
Let $\tau\in(0,1]$, $R\ge2$, $x\in\Lambda_\tau$, $a,b,c\in[0,3]$ with $\Sigma\defeq a+b+c\le6$, and $E_R\defeq\{(u,v)\in\Lambda_\tau^2:|x-u|\le R,|x-v|\le R\}$. Then, with $C$ depending only on an upper bound for $\Sigma$,
\begin{equation}\label{app:a04:eq:PC}
\tau^6\sum_{(u,v)\in E_R}w_a(x-u)w_b(u-v)w_c(v-x)\le C(1+\ln(R/\tau))^4R^{6-\Sigma}.
\end{equation}
Two-point version: $\tau^3\sum_{|x-u|\le R}w_a(x-u)w_c(u-x)\le C(1+\ln(R/\tau))R^{(3-a-c)_+}\tau^{-(a+c-3)_+}$. Tail version: for $s\le3$ and $\lambda>0$, $\tau^3\sum_ue^{-\lambda|x-u|}w_s(x-u)\le C_\lambda(1+\ln\tau^{-1})$.
\end{lemma}
\begin{proof}
\emph{Coincident terms.} $u=x$: $\tau^3w_a(0)=\tau^{3-a}$, and the $v$-sum is $\tau^3\sum_{|v-x|\le R}w_{b+c}(v-x)\le C(1+\ln(R/\tau))R^{(3-b-c)_+}\tau^{-(b+c-3)_+}$ (shells); the product is $\le C(1+\ln(R/\tau))R^{6-\Sigma}$ (if $b+c\ge3$ the $\tau$-exponent is $6-\Sigma\ge0$; if $b+c<3$, $\tau^{3-a}R^{3-b-c}\le R^{6-\Sigma}$). $v=x$ and $u=v$: the same. \emph{Non-coincident terms.} For $r\in\Lambda_\tau\setminus0$ and $r'$ in the sub-cube $Q'_r$ of side $\tau/4$ centred at $r$, $|r'-r|\le\sqrt3\tau/8$, so $|r'|\le2(|r|+\tau)$ and $\tau^3w_s(r)\le64\cdot2^s\int_{Q'_r}|r'|^{-s}dr'$. Applied to the three factors with $u'\in Q'_u$, $v'\in Q'_v$ (so $|(u'-v')-(u-v)|\le\sqrt3\tau/4$ and $|u'-v'|\le2(|u-v|+\tau)$): for distinct lattice points the sub-cubes are disjoint and every pairwise distance among $x,u',v'$ is $\ge\tau-\sqrt3\tau/4\ge\tau/2$, whence
\[
\tau^6\sum_{E_R,\ \text{non-coinc.}}\prod w\le64^2\,2^\Sigma\iint_\Omega|x-u'|^{-a}|u'-v'|^{-b}|v'-x|^{-c}\,du'dv',\quad\Omega\defeq\{|x-u'|,|x-v'|\le R+1,\ \text{all pairwise distances}\ge\tau/2\}.
\]
Dyadic shells in $\varrho\defeq\max(|x-u'|,|x-v'|)\in[\tau/2,R+1]$: scaling $(u'-x,v'-x)=2^i(u'',v'')$ on the shell $2^i\le\varrho<2^{i+1}$ gives $2^{i(6-\Sigma)}J_i$, $J_i\defeq\int_{1\le\varrho''<2,\ \text{distances}\ge\delta_i\defeq\tau/2^{i+1}}|u''|^{-a}|u''-v''|^{-b}|v''|^{-c}$. On the shell one of $|u''|,|v''|$ is $\ge1$, say $|u''|\ge1$, so $|u''|^{-a}\le1$; the remaining $\int_{|v''|\le2}|v''|^{-c}|u''-v''|^{-b}dv''$ with $b,c\le3$ and the two singular points at distance $\ge1$ is $\le C(1+\ln(2^i/\tau))^2$ (a slot equal to $3$ costs one logarithm, a slot $<3$ none), and the $u''$-integral over $1\le|u''|\le2$ is bounded. So $J_i\le C(1+\ln(R/\tau))^2$ and $\sum_{i:\tau/4\le2^i\le R+1}2^{i(6-\Sigma)}\le C(1+\ln(R/\tau))R^{6-\Sigma}$ ($\Sigma\le6$). Total $\le C(1+\ln(R/\tau))^4R^{6-\Sigma}$. Two-point and tail versions: shells in one variable.
\end{proof}

\begin{lemma}[one inner sum]\label{app:a04:lem:Tp}
For $b,c\in[0,3]$ and $u,x\in\Lambda_\tau$: $\tau^3\sum_{v\in\Lambda_\tau,|v-x|\le R}w_b(u-v)w_c(v-x)\le C(1+\ln(R/\tau))\cdot\{w_{b+c-3}(u-x)$ if $b+c>3$; $1$ if $b+c=3$; $R^{3-b-c}$ if $b+c<3\}$.
\end{lemma}
\begin{proof}
Split into $|v-x|\le|u-x|/2$ (there $w_b(u-v)\le2^bw_b(u-x)$), $|v-u|\le|u-x|/2$ (there $w_c(v-x)\le2^cw_c(u-x)$) and the rest (there $|v-u|\ge|u-x|/2$ and $|v-x|\ge|u-x|/2$, so $w_bw_c\le Cw_{b+c}(v-x)$ up to the roles), and sum in shells.
\end{proof}
